\documentclass[11pt]{article}
\usepackage[a4paper,margin=28mm]{geometry}
\usepackage{amsmath,amssymb,amsthm}
\newtheorem{lemma}{Lemma}
\title{Two elementary analytic lemmas for stable polarization}
\author{Independent derivations supporting the power lower-bound proof}
\date{30 September 2026}
\begin{document}\maketitle
These are analytic lemmas, not by themselves a bound on dice sizes.

\begin{lemma}[Centered coefficients from variance alone]
Let $y_1,\ldots,y_N$ be real, $\sum_i y_i=0$, and
$V=\sum_i y_i^2$. For $1\le k\le N$,
\[
 |e_k(y)|\le(eV/k)^{k/2},\qquad
 \frac{|e_k(y)|}{\binom Nk}
       \le\left(\frac{\sqrt{ekV}}{N}\right)^k.
\]
No upper bound on the largest coordinate, and no restriction
$k\le N/2$, is required.
\end{lemma}
\begin{proof}
For every complex $w$,
\[
 |1+w|^2=1+2\operatorname{Re}w+|w|^2
       \le\exp(2\operatorname{Re}w+|w|^2).
\]
Consequently, for every complex $z$,
\[
 \left|\prod_i(1+zy_i)\right|\le\exp(V|z|^2/2).
\]
If $V>0$, Cauchy's coefficient bound on the circle of radius
$\sqrt{k/V}$ gives the first inequality. The case $V=0$ is trivial.
The second follows from $\binom Nk\ge(N/k)^k$.
\end{proof}

\begin{lemma}[Off-diagonal Gaussian bound]
For ordinary Laguerre polynomials $L_j$, $0\le h\le j$, and $x\ge0$,
\[
 x^{h/2}\left|\frac{d^h}{dx^h}L_j(x)\right|
 \le\sqrt{\frac{j!}{(j-h)!}}\,e^{x/2}.
\]
In particular, for $\theta>0$, $z\ge0$, and any finite sum
$Q(z)=\sum_{j=0}^d c_jL_j(\theta z)$ with
$B=\sum_j|c_j|$,
\[
 z^{h/2}|Q^{(h)}(z)|\le B(\theta d)^{h/2}e^{\theta z/2}.
\]
For two such sums $Q_1,Q_2$, both of degree at most $d$,
\[
 |(Q_1Q_2)^{(k)}(z)|
 \le B_1B_2 e^{\theta z}
                \left(2\sqrt{\theta d/z}\right)^k
 \quad(z>0).
\]
\end{lemma}
\begin{proof}
Use the Hilbert space of entire functions with Gaussian norm
\[
 \|f\|^2=\frac1\pi\int_{\mathbb C}|f(w)|^2e^{-|w|^2}\,dA(w).
\]
The monomials $f_j(w)=w^j/\sqrt{j!}$ are orthonormal. For real $a$,
the operator
\[
 (U_af)(w)=e^{-a^2/2+aw}f(w-a)
\]
is unitary, since
$|e^{-a^2/2+aw}|^2e^{-|w|^2}=e^{-|w-a|^2}$.
Expansion of $e^{aw}(w-a)^{j-h}$, followed by extraction of its
$w^j$ coefficient, gives
\begin{align*}
 \langle f_j,U_af_{j-h}\rangle
 &=e^{-a^2/2}\sqrt{\frac{j!}{(j-h)!}}
   \sum_{t=0}^{j-h}\binom{j-h}{t}
        \frac{(-1)^t a^{h+2t}}{(h+t)!}\\
 &=e^{-a^2/2}\sqrt{\frac{(j-h)!}{j!}}
       a^h L_{j-h}^{[h]}(a^2).
\end{align*}
Here square brackets denote the associated-Laguerre parameter,
to distinguish it from differentiation. The coefficient identity is
the defining expansion
$L_{j-h}^{[h]}(x)=\sum_t(-1)^t\binom{j}{h+t}x^t/t!$.
The matrix coefficient has absolute value at most one. Taking
$a=\sqrt{x}$ and using
$D^h L_j=(-1)^hL_{j-h}^{[h]}$ proves the first bound.
The scaled version follows by differentiation, summation, and
$j!/(j-h)!\le d^h$. Derivatives with $h>d$ vanish.
Finally the product rule and the binomial sum $\sum_h\binom kh=2^k$
give the last assertion. At zero, only the weighted or limiting forms
of these statements are intended; no division by zero is used.
\end{proof}
\end{document}
